\section{Introduction}
\label{sec:introduction}

The ability of Large Language Models (LLMs) to engage in self-reflection—identifying and correcting their own errors—is a cornerstone of modern reasoning techniques such as \textsc{Reflexion} \citep{shinn2023reflexion} and \textsc{Self-Refine} \citep{madaan2023selfrefine}. While these methods consistently improve task performance, the underlying mechanism remains a "black box." Does the model truly "re-think" its reasoning, inducing a structural shift in its internal representations, or does it merely re-sample from its output distribution based on a new prompt?

\para{Why this matters.} Understanding whether self-critique reshapes internal states is essential for building more reliable and autonomous AI systems. If reflection is merely shallow re-sampling, the improvements may be brittle and sensitive to prompt phrasing. Conversely, if reflection induces a structured shift toward a more coherent "reasoning subspace," we can leverage these signals for mechanistic control techniques like Representation Engineering (\repe) \citep{yan2025reflctrl} to improve model performance without explicit prompting.

\para{Gap identification.} Recent work has identified specific "reflection vectors" and "self-authorship signals" in the residual stream \citep{zhu2025emergencecontrol,ackerman2024inspection}. However, these studies typically focus on static snapshots of activations. What is missing is a dynamic analysis of how internal states \textit{drift} across multiple rounds of critique and whether this drift correlates with an increasingly structured reasoning manifold.

\para{Our approach.} We propose \methodname, a framework to measure the representational dynamics of LLMs during self-reflection. We extract residual stream activations across multiple rounds of critique and revision on mathematical (\gsm) and commonsense (\csqa) reasoning tasks. We quantify the internal change using cosine distance (drift) and track the linear separability of correct versus incorrect reasoning via logistic regression probes.

\para{Quantitative preview.} Our experiments reveal a stark contrast between autonomous and guided reflection. In the \qwenonefive model, autonomous self-critique leads to representational drift comparable to a simple clarity \rewrite, but fails to improve the linear separability of correct answers (AUC $0.486 \rightarrow 0.444$). However, providing an \external oracle critique causes a massive structural shift, increasing AUC to $0.831$. Surprisingly, the smaller \qwenzerofive model successfully restructures its latent space autonomously (AUC $0.550 \rightarrow 0.719$), highlighting the complex interplay between model scale and reflective capability.

\para{Our main contributions are:}
\begin{itemize}[leftmargin=*,itemsep=0pt,topsep=0pt]
    \item We propose a dynamic framework for measuring \textit{representation drift} and subspace structuring in LLMs during iterative self-reflection.
    \item We demonstrate that representational drift is omnipresent across all revision conditions, but drift alone does not guarantee "cognitive restructuring" or improved separability.
    \item We show that external guidance is significantly more potent at restructuring the latent reasoning space than autonomous self-critique in mid-sized models.
    \item We identify a surprising "scale anomaly" where the smallest model (\qwenzerofive) exhibit more effective autonomous internal restructuring than its larger counterpart (\qwenonefive).
\end{itemize}

The remainder of this paper is organized as follows: \secref{sec:related_work} reviews related work, \secref{sec:methodology} details our experimental setup, \secref{sec:results} presents our findings, \secref{sec:discussion} discusses limitations and implications, and \secref{sec:conclusion} concludes.
